\chapter{认知架构}\label{s:architecture}

我们一直在说心智模型，好像它们是真实存在的东西，
但学习者在学习时，他们的大脑里究竟发生了什么？
简短的回答是我们不知道；
长一点的回答是，我们知道的比过去多得多。
本章会更深入地挖一挖大脑在学习时做些什么，
以及我们如何利用这些来更有效地设计和讲授课程。

\seclbl{那里发生了什么？}{s:architecture-brain}

\figpdf{figures/cognitive-architecture.pdf}{认知架构}{f:arch-model}{}

\figref{f:arch-model}是人的认知架构的一个简化模型。\index{认知架构}
这个模型的核心，
是\secref{s:memory-seven-plus-or-minus}里讨论过的短期记忆与长期记忆之间的分离。
长期记忆就像你家的地下室：\index{长期记忆}
它差不多永久地存放东西，
但你没法直接读取里面的内容。
相反，
你得依靠你的短期记忆，\index{短期记忆}
它是你脑海里那张堆得乱七八糟的厨房桌子。

你需要某样东西时，
大脑把它从长期记忆里取出来，放进短期记忆。
反过来，
进入短期记忆的新信息，
必须经过编码才能存进长期记忆。
如果这些信息没有被编码和存储，
它就不会被记住，学习也就没有发生。

信息主要通过你的言语通道（用于语言）\index{言语通道}
和视觉通道\index{视觉通道}
（用于图像）进入短期记忆\footnote{
  一个更完整的模型
  还会把触觉、嗅觉和味觉也算进去，
  但我们暂且忽略它们。}。
大多数人主要依靠视觉通道，
但当图像和文字彼此互补时，
大脑把两者都记得更好：
它们被一起编码，
所以后来回忆其中之一，有助于触发对另一者的回忆。

语言输入和视觉输入由人脑的不同部位处理，
语言记忆和视觉记忆也是分开存储的。
这意味着，把语言和视觉这两股信息流关联起来需要耗费认知努力：
当某人一边看某段文字、一边听它被念出来时，
他的大脑会不由自主地核对两条通道得到的是不是同样的信息。

所以，当信息以两条不同的通道同时呈现时，学习会增强；
但当这些信息是冗余的而不是互补的时，学习反而会减弱，
这一现象叫做\gref{g:split-attention-effect}{注意分散效应}~\cite{Maye2003}。
例如，
人们通常觉得，从既有旁白、又有屏幕字幕的视频里学习，
比从只有旁白或只有字幕的视频里学习更难，
因为他们得把一部分注意力用在核对外
旁白和字幕是否一致上。
有两个显著的例外：还不太会说这门语言的人，
以及有听觉障碍或其他特殊需求的人；
对这两类人来说，冗余信息带来的价值
可能超过额外的加工负担。

\begin{aside}{一块一块地画}
  注意分散效应解释了为什么
  一边讲一边把图一块一块地画出来，比一次把整张图呈现出来更有效。
  如果图的某些部分和正在说的话同时出现，
  两者就会在学习者的记忆里关联起来。
  之后再指着图上的某一部分，
  就更有可能触发对画那一部分时所说内容的回忆。
\end{aside}

注意分散效应\emph{并不}意味着
学习者不该试着去调和来自多条渠道的信息——毕竟，
这正是他们在真实世界里必须要做的事~\cite{Atki2000}。
它意味着的是：教学不应该要求人们在刚掌握单元技能时就去这么做；
相反，
同时使用多个信息来源，应该被当作一项单独的学习任务。

\begin{aside}{图形并非生来平等}
  \cite{Sung2012}呈现了一项精巧的研究，它区分了\emph{诱惑性}图形\index{图形!诱惑性}
  （非常有趣，但与教学目标并不直接相关）、
  \emph{装饰性}图形\index{图形!装饰性}
  （中性，但与教学目标并不直接相关）
  和\emph{指导性}图形\index{图形!指导性}
  （与教学目标直接相关）。
  拿到任何一类图形的学习者，对材料的满意度评分都比没拿到图形的学习者高，
  但只有拿到指导性图形的学习者，实际表现才更好。

  同样，~\cite{Stam2013,Stam2014}发现，
  信息更多反而可能让表现更差。
  他们向孩子们呈现图片、图片加数字，或只有数字，来完成两项任务。
  对一些人来说，
  有图片或有图片加数字的表现胜过只有数字，
  但对另一些人来说，
  有图片胜过图片加数字，而图片加数字又胜过只有数字。
\end{aside}

\seclbl{认知负荷}{s:architecture-load}

在~\cite{Kirs2006}中，Kirschner、Sweller 和 Clark 写道：

\begin{quote}

  尽管无引导或最少引导的教学方法
  非常流行、也符合直觉{\ldots}这些方法忽略了
  构成人类认知架构的那些结构，
  也忽略了半个世纪以来实证研究得到的证据，
  这些证据一致表明，最少引导的教学在效果和效率上
  都低于那些强烈强调引导学生学习过程的教学方法。
  只有当学习者已经拥有足够高的先备知识、能提供「内部」引导时，
  引导的优势才开始消退。

\end{quote}

抛开那些行话，
作者们的主张是：让学习者自己提问题、
自己定目标、
自己找一条穿过某个学科的路径，
不如一步步教他们怎么做有效。
这种「自己选冒险路线」的做法叫做\gref{g:inquiry-based-learning}{探究式学习}，
它很符合直觉：
毕竟，
谁会\emph{反对}让学习者主动去以贴近现实的方式解决真实世界的问题呢？
然而，
在一个全新的领域里要求学习者这么做，会把他们压垮，
因为这要求他们同时掌握一个领域的事实内容
和它的问题解决策略。

更具体地说，
\gref{g:cognitive-load}{认知负荷理论}提出，
人们在学习时必须应付三样东西：

\begin{description}

\item[\grefdex{g:intrinsic-load}{内在负荷}{认知负荷!内在}]
  是人们为了吸收新材料而必须放在心上的东西。

\item[\grefdex{g:germane-load}{关联负荷}{认知负荷!关联}]
  是把新信息与旧信息联系起来所需的（有益的）心理努力，
  这也是区分「学习」与「记忆」的要素之一。

\item[\grefdex{g:extraneous-load}{外在负荷}{认知负荷!外在}]
  是任何会分散学习注意力的东西。

\end{description}

认知负荷理论认为，
人们必须在一份固定的工作记忆里，把它分配给这三样东西。
作为老师，我们的目标是让可用于处理内在负荷的记忆最大化，
这就意味着在每一步都减少关联负荷、并消除外在负荷。

\subsection*{帕森斯问题}

有一类可以用认知负荷来解释的练习，
在语言教学中经常使用。
假设你让某人把「How is her knee today?」这句话翻译成弗里西语。
要解决这个问题，
他们既得回忆起词汇、又得回忆起语法，
这是双重的认知负荷。
但如果你让他们把「hoe」「har」「is」「hjoed」「knie」这几个词排成正确的顺序，
你就是在让他们把注意力完全集中在学语法上。
不过，如果你把这几个词写成五种不同的字体或颜色，
你就增加了外在认知负荷，
因为他们会不由自主地（可能还不知不觉地）花一些力气，
去弄清这些差异有没有意义
（\figref{f:architecture-frisian}）。

\figimg{figures/frisian.png}{构造一个句子}{f:architecture-frisian}{}

这件事在编程里的对应物，
叫做\gref{g:parsons-problem}{帕森斯问题}\footnote{以它的创建者之一命名。}~\cite{Pars2006}。
教人编程时，
你可以把解决问题需要的几行代码给他们，
让他们排出正确的顺序。
这样他们就能专注于控制流和数据依赖，
而不被变量命名分心，也不必费劲去记该调用哪些函数。
多项研究表明，帕森斯问题让学习者花的时间更少，
却产生同等的教学效果~\cite{Eric2017}。

\subsection*{渐隐示例}

另一类可以用认知负荷来解释的练习，
是给学习者一系列\grefdex{g:faded-example}{渐隐示例}{渐隐示例}。
一系列中的第一个例子，完整地展示某种问题解决策略的用法。
下一个问题类型相同，
但留了几个空让学习者填。
每往后一个问题，给学习者的\gref{g:scaffolding}{脚手架}都更少，
直到最后要求他们从头独立解决一个完整的问题。
例如，在教高中数学代数时，
我们可能先给出这个：

\begin{center}
\begin{tabular}{rcl}
  (4x + 8)/2	& = &	5	\\
  4x + 8	& = &	2 * 5	\\
  4x + 8	& = &	10	\\
  4x		& = &	10 - 8	\\
  4x		& = &	2	\\
  x		& = &	2 / 4	\\
  x		& = &	1 / 2
\end{tabular}
\end{center}

\noindent
然后让学习者解这个：

\begin{center}
\begin{tabular}{rcl}
  (3x - 1)*3	& = &	12	\\
  3x - 1	& = &	\_ / \_	\\
  3x - 1	& = &	4	\\
  3x		& = &	\_	\\
  x		& = &	\_ / 3	\\
  x		& = &	\_
\end{tabular}
\end{center}

\noindent
再解这个：

\begin{center}
\begin{tabular}{rcl}
  (5x + 1)*3	& = &	4	\\
  5x + 1	& = &	\_ 	\\
  5x		& = &	\_ 	\\
  x		& = &	\_
\end{tabular}
\end{center}

\noindent
最后是这个：

\begin{center}
\begin{tabular}{rcl}
  (2x + 8)/4	& = &	1	\\
   x		& = &	\_
\end{tabular}
\end{center}

在教 Python 时，一个类似的练习可以先给学习者看\index{Python}
如何求一个单词列表的总长度：

\begin{minted}{text}
# total_length(["red", "green", "blue"]) => 12
define total_length(list_of_words):
    total = 0
    for word in list_of_words:
        total = total + length(word)
    return total
\end{minted}

\noindent
然后让他们填出下面这段的空
（这会把他们的注意力集中在控制结构上）：

\begin{minted}{text}
# word_lengths(["red", "green", "blue"]) => [3, 5, 4]
define word_lengths(list_of_words):
    list_of_lengths = []
    for ____ in ____:
        append(list_of_lengths, ____)
    return list_of_lengths
\end{minted}

下一个问题可能是这个
（这会把他们的注意力集中在更新最终结果上）：

\begin{minted}{text}
# join_all(["red", "green", "blue"]) => "redgreenblue"
define join_all(list_of_words):
    joined_words = ____
    for ____ in ____:
        ____
    return joined_words
\end{minted}

最后会要求学习者自己写出一个完整的函数：

\begin{minted}{text}
# make_acronym(["red", "green", "blue"]) => "RGB"
define make_acronym(list_of_words):
    ____
\end{minted}

渐隐示例之所以有效，
是因为它把问题解决策略一块一块地引入：
每一步，
学习者只需要应付一个新问题，
这比面对一块空屏幕或一张白纸更不吓人（\secref{s:classroom-practices}）。
它也鼓励学习者去思考各种方法之间的异同，
这有助于在他们的心智模型里建立连接，从而帮助回忆。

构造一个好的渐隐示例，关键在于
想清楚它打算教的是什么问题解决策略。
例如，
上面那些编程问题全都用了累加器设计模式，
也就是把处理集合中各个条目的结果
以某种方式反复加到同一个变量上，从而得出最终结果。

\begin{aside}{认知学徒制}
  另一种同样使用脚手架和渐隐的学习与教学模型，
  是\gref{g:cognitive-apprenticeship}{认知学徒制}，
  它强调师傅把技能和洞见传授给学徒的方式。
  师傅提供表现和结果的范本，
  然后通过解释自己在做什么、为什么这样做来指导新手~\cite{Coll1991,Casp2007}。
  学徒反思自己的问题解决过程，
  例如把思路说出来，或评论自己的作品，
  最终去探索自己选择的问题。

  这个模型告诉我们：
  老师在呈现一个新想法时应该给出好几个例子，
  好让学习者看出该概括什么；
  我们还应该变换问题的形式，以讲清
  哪些是表面特征、哪些不是\footnote{有很长一段时间，
    我相信用来存放函数将要返回的值的那个变量
    \emph{必须}叫做 \texttt{result}，
    因为我的老师在例子里总是用这个名字。}。
  问题应该放在真实世界的语境里呈现，
  我们还应该鼓励自我解释，帮学习者组织和弄懂刚学到的内容
  （\secref{s:individual-strategies}）。
\end{aside}

\subsection*{标注子目标}

\grefdex{g:subgoal-labeling}{标注子目标}{标注子目标}是指
为问题解决过程的分步描述中的各个步骤命名。
\cite{Marg2016,Morr2016}发现，有标注子目标的学习者
在解帕森斯问题上比没有标注的学习者做得更好，
其他领域也观察到同样的好处~\cite{Marg2012}。
回到前面用过的 Python 例子，
在求单词列表总长度或构造首字母缩写时，子目标是：

\begin{enumerate}

\item
  创建一个要返回类型的空值。

\item
  从循环变量中取得要加到结果上的值。

\item
  用那个值更新结果。

\end{enumerate}

标注子目标之所以有效，是因为把相关步骤归并成有名字的组块（\secref{s:memory-seven-plus-or-minus}）\index{组块}
有助于学习者区分，什么是通用的、什么是当前问题特有的。
它还有助于他们为那类问题建立一个心智模型，
从而能解决同类别的其他问题，
并给他们一个自我解释的自然机会（\secref{s:individual-strategies}）。

\subsection*{最简手册}

认知负荷理论最纯粹的应用，也许是 John Carroll 的\index{Carroll, John}
\gref{g:minimal-manual}{最简手册}~\cite{Carr1987,Carr2014}。
它的出发点是引用一位用户的话：
「我想做一件事，而不是学会怎么做所有的事。」
Carroll 和同事重新设计了培训，把每个想法都呈现为一页独立的自包含任务：
一个描述这一页讲什么的小标题、
只讲一件事怎么做的分步说明
（例如，如何在一个文本编辑器里删除一个空行），
然后是若干条关于如何识别和调试常见问题的说明。
他们发现，以这种方式重写培训材料，总体上让材料变得更短，
而且使用这些材料的人学得更快。
后续研究证实，无论使用者此前有没有计算机经验，
这种方法都胜过传统方法~\cite{Lazo1993}。
\cite{Carr2014}这样总结这项工作：

\begin{quote}

  我们的「极简主义」设计力图利用用户的主动性和先备知识，
  而不是用警告和规定的步骤去控制它。
  它强调，用户通常会给这种学习带来大量专长和洞见，
  例如关于任务领域的知识，
  并且这类知识可以成为教学设计者的资源。
  极简主义利用错误识别、诊断和恢复的环节，
  而不是仅仅试图预防错误。
  它把排错和恢复框定为学习机会，而不是异常。

\end{quote}

\seclbl{其他的学习模型}{s:architecture-theory}

认知负荷理论的批评者有时主张，\index{认知负荷!批评}
任何结果都可以事后通过把损害表现的东西贴上外在负荷的标签、
把无害的东西贴上内在或关联负荷的标签来合理化。
然而，
基于认知负荷理论的教学是不可否认地有效的。
例如，
\cite{Maso2016}重新设计了一门数据库课程，以消除注意分散和冗余效应，
并提供示例和子目标。
新课程把考试不及格率降低了 34\%，
还提高了学习者满意度。

在~\cite{Kirs2006}发表十年之后，
越来越多的人相信，如果以正确的方式看待，
认知负荷理论和探究式方法是可以兼容的。
\cite{Kaly2015}主张，认知负荷理论基本上是在一个更大的语境里对学习进行微观管理，
这个语境还考虑动机之类的东西；
而~\cite{Kirs2018}把认知负荷理论扩展到包含学习的协作层面。
正如~\cite{Mark2018}（在\secref{s:individual-strategies}里讨论过）那样，
研究者的视角可能不同，
但他们理论的实际落地方式往往殊途同归。

教育研究中的一个挑战是，
一旦你跳出标准化的西方课堂，
我们所说的「学习」究竟是什么，就变得很复杂。
\gref{g:educational-psychology}{教育心理学}里有两个具体视角影响了本书。
我们到目前为止用的是\gref{g:cognitivism}{认知主义}，
它关注模式识别、记忆形成和回忆之类的东西。
它擅长回答低层次的问题，
但通常忽略更大的问题，比如
「我们所说的『学习』是什么意思？」
以及「谁有资格来定义它？」
另一个是\gref{g:situated-learning}{情境学习}，
它关注把人带进一个共同体，
并承认
教与学总是植根于我们是谁、以及我们渴望成为谁。
我们会在\chapref{s:community}里更详细地讨论它。

\hreffoot{http://www.learning-theories.com/}{Learning Theories 网站}
和~\cite{Wibu2016}
对这些以及其他视角都有很好的概述。
除认知主义之外，
最常遇到的还包括\gref{g:behaviorism}{行为主义}
（把教育当作刺激/反应的条件作用）、
\gref{g:constructivism}{建构主义}
（把学习看作一个主动的过程，学习者在此过程中为自己建构知识）、
以及\gref{g:connectivism}{联通主义}
（认为知识是分布式的，
学习就是建立、生长和修剪连接的过程，
并强调互联网使学习具有的社会性）。
这些视角能帮我们整理思路，
但在实践中，
我们总得在课堂上、面对真实的学习者去尝试新方法，
才能弄清它们在多大程度上摆平了各种相互作用的力量。

\seclbl{练习}{s:architecture-exercises}

\exercise{创建一个渐隐示例}{两人一组}{30}

程序里很常见的一种操作，是统计有多少东西分属不同类别：
例如，
一个图像里不同颜色各出现了多少次，
或者一段文字里不同的词各出现了多少次。

\begin{enumerate}
\item
  创建一个简短的例子（不超过 10 行代码），向人们展示怎么做这件事，
  然后再创建第二个例子，以相似的方式解决一个相似的问题，
  但留出几个空让学习者填。
  你是怎么决定哪些地方该渐隐掉的？
  这一系列里的下一个例子会是什么？

\item
  为你的例子确定受众。
  例如，
  这些是只懂一些基础编程概念的初学者吗？
  还是已经有一些编程经验的学习者？

\item
  把你的例子给同伴看，
  但\emph{不要}告诉他们你觉得这是给哪个水平的。
  等他们填完空之后，
  让他们猜一猜目标水平。
\end{enumerate}

如果参加培训的人里有完全不会编程的，
尽量把他们分到不同的小组里，
让他们在那些小组里扮演学习者。
或者，
换一个别的任务领域，为它做一个渐隐示例。

\exercise{给负荷分类}{小组}{15}

\begin{enumerate}

\item
  选一节你们小组某个成员最近教过或上过的短课。

\item
  把其中包含的想法、指令和解释列成一个要点清单。

\item
  把它们分别归为内在负荷、关联负荷或外在负荷。
  你们一致同意的是哪些？
  在哪里有分歧？为什么？

\end{enumerate}

（\secref{s:memory-exercises}里的练习「注意你的盲区」
会让你明白这个要点清单该做到多细。）

\exercise{创建一个帕森斯问题}{两人一组}{20}

写五六行能做某件有用之事的代码，
把它们打乱，
让你的同伴排出顺序。
如果你用的是像 Python 这样靠缩进的语言，
就不要给任何行加缩进；
如果你用的是像 Java 这样用花括号的语言，
就不要包含任何花括号。
（如果你的小组里有不会编程的人，
换一个别的任务领域，比如做香蕉面包。）

\exercise{最简手册}{个人}{20}

写一页指南，讲怎么完成学习者可能在你课上遇到的某件事，
比如把文字水平居中，
或者打印一个带指定位数小数的数字。
尽量列出学习者可能看到的至少三四种错误行为或结果，
并为每一种加一两行说明，
解释它为什么会发生、该怎么纠正。

\exercise{认知学徒制}{两人一组}{15}

挑一个你两三分钟就能做完的编程问题，
一边做一边把思路说出来，
而你的同伴就你在做什么、为什么这么做提问。
不要只解释你在做什么，
还要解释你为什么这么做、
你怎么知道这是对的做法、
以及你考虑过但又放弃了哪些别的做法。
做完之后，
和同伴交换角色，重复这个练习。

\exercise{示例}{两人一组}{15}

看例子能让人比只写大量代码更快地学会编程~\cite{Skud2014}，
而通过追踪或调试来拆解代码也能增加学习~\cite{Grif2016}。
两人一组，
拿一段 10—15 行的代码，解释每一条语句做什么、
以及为什么需要它。
花了多长时间？
你觉得每行代码需要解释的东西有多少？

\exercise{评论图形}{个人}{30}

\cite{Maye2009,Mill2016a}提出了好教学图形的六条原则：

\begin{description}

\item[信号提示：]
  在视觉上突出最重要的点，
  让它们从不太关键的材料中脱颖而出。

\item[空间邻近：]
  把图注放得尽量靠近图形，以抵消在两者之间来回切换的代价。

\item[时间邻近：]
  让口头旁白和图形在时间上尽量靠近。
  （两者同时呈现比一前一后呈现更好。）

\item[分段：]
  在呈现一长串材料，或学习者对主题没有经验时，
  把呈现切分成短小的片段，
  让学习者自己控制以多快进入下一段。

\item[预训练：]
  如果学习者不知道你的呈现里用到的主要概念和术语，
  就在事先只教这些概念和术语。

\item[模态：]
  人们从图片加旁白中学得比从图片加文字中学得好，
  除非他们是非母语者，
  或者其中含有技术词汇或符号。

\end{description}

在网上选一段使用幻灯片或其他静态呈现的课程或演讲视频，
按这六条标准，把它的图形评为「差」「中」或「好」。

\section*{回顾}

\figpdfhere{figures/conceptmap-cognitive-load.pdf}{概念：认知负荷}{f:architecture-cognitive-load}{}
